For Client and Server key/certificate generation I used the program \lstinline{openssl}\ref{openssl.keygen}.

The following code snippet showcases the \textit{Server} key generation of \lstinline{00_keygen.sh} script. Please note: For the actual test environment, the \lstinline{00_keygen.sh} should be sourced\cite{shell.bash} within the container environment,as can be seen within the \textit{Dockerfiles}\ref{apx:dockerfile.serverfile, apx:dockerfile.clientfile}

\subsection{Overview of openssl key generation with \textit{00\_keygen.sh}}
\label{apx:lab-final_keygen}

\begin{lstlisting}[language=bash, caption=openssl keygeneration script]
#!/bin/ash

# ========================================
# Openssl key/certificate generator script
# ----------------------------------------

# Host Key/Certification generation according to `cert.conf` contents.
# The functions accepts 2 input parameters:
# 	`server`	:= server key/certificate generation
# 	`client` 	:= client key/certificate generation 

# Please note: For this assignment I vpn-server is the same as CA and Intermediary Server

# The file is being sourced trough dotfiles, hence
# the function is usable as a command as follows:
#
# host_keygen server	:= generates server ca/key
# host_keygen client	:= generates client ca/key

# Please note: The script is meant to be used as beign sourced
# Please refer to the `source` command arguments in the dockerfiles
host_keygen() {

echo
echo "======================================================"
echo "Client/Server keypair generation using OpenSSL library"
echo "======================================================"
echo

# Choose Server/Client
HOST_MODE=$1
	case $HOST_MODE in
		
		# Generating server key/certificate using group 'prime256v1' according to RFC 3279
		server)
	
		export NAME_TAG="$(hostname)"
		export CERT_CONF="cert.conf"
		export CA_KEY_NAME="CA_KeyFile"
		export CA_CERT_NAME="CA_Certificate"
		export TRUST_STORE="/etc/ssl/certs/"
		export TRUST_KEY_STORE="/etc/ssl/private/"
		export SHARED_DIR="/volumes/inter"
			echo
			echo "Generating the Root CA certificate as $NAME_TAG"
			echo "---------------------------------------------------------"
			echo
			
				# Key/CA generation based on 2048-bit RSA
				# .......................................
				openssl genrsa -out "$CA_KEY_NAME".pem 2048
				
				echo "Generate the Root CA's self-signed certificate: "
				echo "---------------------------------------------------------"
				openssl req -new -x509 -days 365 -sha256 -key "$CA_KEY_NAME".pem -config "$CERT_CONF" -out "$CA_CERT_NAME".crt

				echo 'Check Certificate Content:'
				echo "---------------------------------------------------------"
				openssl x509 -in "$CA_CERT_NAME".crt -text -noout

				# Key/CA generation based on Elliptice Curves
				# ...........................................
				# Store keyfiles un-encrypted on disk (PEM):
				#openssl ecparam -name prime256v1 -genkey -noout -out "$NAME_TAG"_ec_server.pem
                # Generate a Certificate Signing Request (CSR)
				# Template: cert.conf
				#openssl req -new -key "$NAME_TAG"_ec_server.pem -out "$NAME_TAG"_ec_server.csr --config ca.conf

		export INTER_KEY_NAME="Inter_KeyFile"
		export INTER_CSR_NAME="Inter_CSR_Certificate"
		export INTER_CERT_NAME="Inter_Certificate"
		export COMBINED_CERT_NAME="combined"
			echo	
			echo "Generating Intermediary certificate as $NAME_TAG"
			echo "---------------------------------------------------------"
			echo
				openssl genrsa -out "$INTER_KEY_NAME".pem 2048
				openssl req -new -sha256 -key "$INTER_KEY_NAME".pem -config "$CERT_CONF" -out "$INTER_CSR_NAME".csr

			echo 
			echo 'Issue the signed certificate for Intermediate Certificate'
			echo "---------------------------------------------------------"
			openssl x509 -days 365 -req -in "$INTER_CSR_NAME".csr -CA "$CA_CERT_NAME".crt -CAkey "$CA_KEY_NAME".pem -CAcreateserial -out "$INTER_CERT_NAME".crt
			
			echo
			echo "Updating Trusted Stores of Server: $NAME_TAG"
			echo "---------------------------------------------------------"
			
			# Create a combined cert: 
			cat "$INTER_CERT_NAME".crt "$CA_CERT_NAME".crt > "$COMBINED_CERT_NAME".crt
			
			cp "$CA_CERT_NAME".* "$INTER_CERT_NAME".crt "$COMBINED_CERT_NAME".crt "$TRUST_STORE"
			cp *.pem "$TRUST_KEY_STORE"
			update-ca-certificates
			echo "Trusted Store updated."

			echo
			echo "Moving files to $SHARED_DIR"
			echo "---------------------------------------------------------"
			cp "$CA_CERT_NAME".* "$INTER_CERT_NAME".crt "$COMBINED_CERT_NAME".crt "$SHARED_DIR"
						
			
				# Key/CA generation based on Elliptice Curves
				# ...........................................
				# Generate a Self-Signed Server Certificate:
				#openssl x509 -req -in "$NAME_TAG"_ec_server.csr -signkey "$NAME_TAG"_ec_server.pem -out "$NAME_TAG"_ec_server-cert.pem -days 365
                # Check the Server certification:
				#openssl x509 -in "$NAME_TAG"_ec_server-cert.pem -text -noout
                # Check the private key
				#openssl ec -in "$NAME_TAG"_ec_server.pem -text -noout
		
			#echo
			#echo "Writing file checksums to SHA256SUMS_"$NAME_TAG"_ec_server.txt"
			#echo ".............................................................."
			#echo

			# Writing sha256sums of the generated files
			#sha256sum "$NAME_TAG"_ec_server* > SHA256SUMS_"$NAME_TAG".txt
		unset NAME_TAG
		unset CERT_CONF
		unset CA_KEY_NAME
		unset CA_CERT_NAME
		unset INTER_KEY_NAME
		unset INTER_CSR_NAME
		unset INTER_CERT_NAME
		unset SHARED_DIR
			;;
		
		# Generating client key/certificate using group "prime256v1" according to RFC 3279
		client)

		export NAME_TAG="$(hostname)"
		export CERT_CONF="cert.conf"
		export CA_KEY_NAME="CA_KeyFile"
		export CA_CERT_NAME="CA_Certificate"
		export INTER_KEY_NAME="Inter_KeyFile"
		export INTER_CSR_NAME="Inter_CSR_Certificate"
		export INTER_CERT_NAME="Inter_Certificate"

		export USER_KEY_NAME="User_KeyFile"
		export USER_CSR_NAME="User_CSR_Certificate"
		export USER_CERT_NAME="User_Certificate"
		export USER_PUBKEY_NAME="PublicKey"

		export COMBINED_CERT_NAME="combined"
		export TRUST_STORE="/etc/ssl/certs/"
		export TRUST_KEY_STORE="/etc/ssl/private/"
		export SHARED_DIR="/volumes/inter"


			echo
			echo "Generating Client Certificate for $NAME_TAG"
			echo "---------------------------------------------------------"
			echo
			
			# Generate User's RSA key.
			openssl genrsa -out User_KeyFile.pem 2048
			
			# Generate a CSR that includes the intermediate certificate public key.
			# Please note, for this assignment CA is equalst to Intermediary Server
			openssl req -new -sha256 -key "$USER_KEY_NAME".pem -config "$CERT_CONF" -out "$USER_CSR_NAME".csr
			
			# Allow the intermediate certificate to issue the signed certificate for the user certificate.
			openssl x509 -days 365 -req -in "$USER_CSR_NAME".csr -CA "$INTER_CERT_NAME".crt -CAkey "$INTER_KEY_NAME".pem -CAcreateserial -out "$USER_CERT_NAME".crt
			
			# Get User public key
			openssl x509 -pubkey -noout -in "$USER_CERT_NAME".crt > "$USER_PUBKEY_NAME".pem
			
			echo
			echo "Updating Trusted Stores of Client: $NAME_TAG"
			echo "---------------------------------------------------------"
			
			cp /volume/inter/* .
			cp "$CA_CERT_NAME".* "$INTER_CERT_NAME"* "$COMBINED_CERT_NAME".crt   "$USER_CERT_NAME"* "$TRUST_STORE"
			cp *.pem "$TRUST_KEY_STORE"
			
			update-ca-certificates
			echo "Trusted Store updated."
				# Key/CA generation based on Elliptice Curves
				# ...........................................
				# Store keyfiles un-encrypted on disk (PEM):
                #openssl ecparam -name prime256v1 -genkey -noout -out "$NAME_TAG"_ec_client.pem
                # Generate a Certificate Signing Request (CSR)
                # Template: ca.conf
                #openssl req -new -key "$NAME_TAG"_ec_client.pem -out "$NAME_TAG"_ec_client.csr --config ca.conf



				# Key/CA generation based on Elliptice Curves
				# ...........................................
				# Generate a Self-Signed Client Certificate:
                #openssl x509 -req -in "$NAME_TAG"_ec_client.csr -signkey "$NAME_TAG"_ec_client.pem -out "$NAME_TAG"_ec_client-cert.pem -days 365
                # Check the Client certification:
                #openssl x509 -in "$NAME_TAG"_ec_client-cert.pem -text -noout
                # Check the private key
                #openssl ec -in "$NAME_TAG"_ec_client.pem -text -noout

#			echo
#			echo "Writing file checksums to SHA256SUMS_"$NAME_TAG"_ec_client.txt"
#			echo ".............................................................."
#			echo

			# Writing sha256sums of the generated files
#			sha256sum "$NAME_TAG"_ec_client* > SHA256SUMS_"$NAME_TAG".txt

		unset NAME_TAG


			;;

		*)
			echo "Error: Key/ca generation failed. Argument is missing: Write 'server' or 'client' as an argument !"
			echo
			echo "Example usage for Server key/certificate generation: "
			echo "host_keygen server"
			echo
			echo "Example usage for Client key/certificate generation:"
			echo "host_keygen client"
			;;
esac
}
\end{lstlisting}



The \lstinline{openssl} script uses the config file \lstinline{cert.conf} for certificate generation. The following snippet showcases a small piece of the code, for the full content please refer to the file \lstinline{cert.conf}: 

\begin{lstlisting}[language=bash, caption=config file snippet for openssl certification template ]
[ req ]
default_bits       = 2048
default_md         = sha256
distinguished_name = req_distinguished_name
prompt             = no
string_mask        = utf8only
req_extensions     = req_ext

[ req_distinguished_name ]
C            = DA
ST           = AARHUS
L            = AARHUS
O            = Aarhus_University
OU           = Systems_Security
CN           = vpn-server-10.9.0.11
emailAddress = @post.au.dk

[ req_ext ]
subjectAltName = @alt_names

[ alt_names ]
DNS.1 = vpn-server-10.9.0.11
IP.1  = 10.9.0.11

[ v3_ca ]
subjectKeyIdentifier   = hash
authorityKeyIdentifier = keyid:always,issuer
basicConstraints       = critical, CA:true
keyUsage               = critical, digitalSignature, cRLSign, keyCertSign

[ v3_intermediate_ca ]
subjectKeyIdentifier   = hash
authorityKeyIdentifier = keyid:always,issuer
basicConstraints       = critical, CA:true, pathlen:0
keyUsage               = critical, digitalSignature, cRLSign, keyCertSign

[ usr_cert ]
basicConstraints       = CA:FALSE
nsCertType             = client, email
nsComment              = "Client Certificate"
subjectKeyIdentifier   = hash
authorityKeyIdentifier = keyid,issuer
keyUsage               = critical, digitalSignature, keyEncipherment
extendedKeyUsage       = clientAuth, emailProtection

\end{lstlisting}

For the complete generation of Client/Server side key/certification generation
please see  the file \lstinline{00_keygen.sh}.

